\chapter*{Internship Logbook}
\addcontentsline{toc}{chapter}{Internship Logbook}

\section*{Week 1}
I started my internship by understanding the operational context at Africom Services and mapping the real intervention workflow from request planning to client reporting. My main tasks were to read the PRD in depth, discuss business constraints with my supervisor, and compare the target digital process with the existing fragmented process used by field teams. I learned that the biggest challenge was not only technical implementation but also process standardization between coordinators, supervisors, and technicians. I used requirement analysis skills, system thinking, and technical documentation reading to identify critical pain points such as delayed validation, inconsistent field data, and weak traceability. The main difficulty I faced was translating broad business statements into concrete system behaviors; I solved this by rewriting each business objective as a measurable feature and linking it to actors and expected outputs. This first week confirmed that the internship matched my expectation of full-stack engineering with strong business impact, and it motivated me to focus on building a solution that is both technically robust and operationally useful.

\section*{Week 2}
During the second week, I transformed the initial understanding into a structured planning approach using a SCRUM cadence adapted to a single-developer context. I defined sprint boundaries, backlog themes, and acceptance criteria, while selecting the final stack (Spring Boot, Angular, Flutter, PostgreSQL, Firebase, Apache POI). I learned how important it is to align architecture decisions with non-functional requirements from the beginning, especially security, maintainability, and performance. I used planning, prioritization, and architectural design skills to break the project into five coherent sprints, each with clear objectives and deliverables. A challenge appeared when balancing ambition and timeline: many features were interconnected, and implementing one too early would create rework. I addressed this by sequencing work from foundation to business features, then quality governance, then reporting and optimization. My reflection this week is that disciplined planning reduces stress later and makes technical trade-offs much clearer.

\section*{Week 3}
This week focused on formal requirements engineering and traceability. I produced actor-role mappings, functional requirement catalogs, non-functional targets, and API-oriented scope definitions that later guided development. I learned to convert product language into testable requirements, for example turning the need for better supervision into explicit endpoints for pending validation queues, rejection feedback, and status history. I used UML thinking, API contract design, and domain modeling to define entities like RapportRequest, RapportTemplate, Files, ExecutionDetails, StatusHistory, and ConfigExport. The main difficulty was handling dynamic data requirements for field forms without creating a rigid schema. I solved this by planning JSONB-based structures in PostgreSQL for template mapping and report line data, while keeping critical governance data relational. This week improved my ability to bridge business and engineering, and I realized that good requirement quality directly influences code quality.

\section*{Week 4}
I launched Sprint 1 by setting up the backend, web, and mobile foundations and establishing secure project workspaces. My tasks included Spring Boot initialization, security configuration scaffolding, baseline database schema setup, and preparing Angular and Flutter authentication flows. I learned practical lessons about early infrastructure choices, especially how authentication and error handling conventions influence every later feature. I used Java/Spring configuration skills, Git workflow discipline, and environment management to keep components aligned across platforms. A challenge emerged around synchronizing token-based authentication behavior between web and mobile clients, because session handling patterns differ. I solved this by standardizing the authentication contract and defining consistent token propagation and validation behavior from day one. My reflection is that investing effort in clean foundations saved considerable debugging time in later weeks.

\section*{Week 5}
In the second week of Sprint 1, I implemented JWT authentication and role-based access control with Spring Security and endpoint-level authorization. I developed core identity endpoints such as /api/authenticate, /api/account, /api/account/change-password, and baseline administration routes for users. I learned how role boundaries affect business trust, since each persona (technician, supervisor, coordinator, admin) needs a different and safe access surface. I used backend security engineering skills, validation logic, and exception handling to produce predictable API behavior with standardized error payloads. The challenge was ensuring that unauthorized access attempts were not only blocked but also clearly diagnosable during testing. I solved this by adding structured logging and clear response semantics for forbidden and unauthorized scenarios. This week strengthened my confidence in secure API design and reminded me that correctness includes both enforcement and observability.

\section*{Week 6}
I completed Sprint 1 by integrating Angular and Flutter clients with the authentication layer and validating end-to-end login and session restoration scenarios. On Angular, I implemented route guards, token interceptors, and role-sensitive navigation. On Flutter, I implemented BLoC-driven login state transitions, secure token persistence, startup session restoration, and controlled logout. I learned that cross-platform consistency is essential for user trust, especially when users switch between web supervision and mobile execution contexts. I used TypeScript reactive forms, Dart BLoC patterns, and API integration testing skills. The main challenge was handling edge cases such as expired tokens and partially restored sessions. I solved this through explicit client-side state resets and centralized unauthorized flow handling. My reflection is that user-facing reliability often depends on invisible details in session management.

\section*{Week 7}
I began Sprint 2 by implementing intervention request creation and assignment workflows. My tasks included building RapportRequest domain logic, creating request endpoints, and ensuring initial lifecycle states were correctly applied (starting from PLANNED). I learned how workflow modeling must reflect operational responsibility, because assignment data, planning metadata, and ownership rules drive all downstream actions. I used Spring service-layer design, DTO mapping, and PostgreSQL relational modeling skills to ensure data consistency and traceability. A challenge appeared when validating mandatory metadata without over-restricting business users. I solved it by applying layered validation: structural checks at DTO level and contextual checks in service logic. This week showed me that business-critical workflows need both strict control and practical flexibility.

\section*{Week 8}
This week I delivered template management and dynamic form governance. I implemented CRUD for RapportTemplate and introduced mapping JSON handling so forms could be configured without code changes for every intervention type. I learned that dynamic systems require careful contract design; a small schema ambiguity can break rendering or validation on mobile. I used JSON schema thinking, backend validation strategies, and frontend form architecture skills to keep template definitions robust. The challenge was preventing invalid template configurations from propagating to execution screens. I solved it by adding consistency checks before template activation and by testing sample mappings against realistic field scenarios. My reflection is that configurability is powerful, but it only works when governance rules are explicit.

\section*{Week 9}
I finalized Sprint 2 with filtering, lookup, and data-quality enforcement across web and mobile experiences. I implemented paginated search by status/client/reference in Angular and ensured mobile dynamic forms respected required and typed constraints before submission. I learned that usability and data quality are tightly connected: better filtering and clearer validation messages reduce operational friction. I used Angular component design, API query parameter handling, and validation UX skills. The main difficulty was keeping response consistency between list views and detail views while adding multiple filter combinations. I solved this by standardizing query specifications and validating edge cases with Postman scenarios. This week reinforced my understanding that small UI and API consistency decisions have large effects on daily operations.

\section*{Week 10}
I started Sprint 3 by implementing file management in the backend, including secure upload, metadata persistence, and intervention linkage. I created and integrated file-related entities and services, including upload constraints and image-specific processing logic. I learned practical storage trade-offs, especially balancing evidence quality with storage and transfer costs. I used Spring multipart handling, service orchestration, and metadata modeling skills. The major challenge was managing large files and avoiding uncontrolled storage growth. I solved this by enforcing file size/type validation and integrating image compression in the upload pipeline. My reflection is that evidence management is not just attachment storage; it is a core trust mechanism for supervisory validation.

\section*{Week 11}
This week I focused on mobile evidence capture and geolocation tracking. I implemented Flutter BLoC flows for file upload, offline queuing when connectivity fails, and GPS capture with permission handling. I learned that field conditions impose different design priorities than office software: resilience and offline continuity become mandatory. I used Dart asynchronous programming, local queue design, and geolocation service integration skills. The challenge was handling unreliable networks without frustrating technicians or losing data integrity. I solved this by queueing failed uploads locally and syncing later, while giving users clear feedback about upload status. This week changed my perspective on mobile development by showing how strongly context shapes architecture.

\section*{Week 12}
I completed Sprint 3 and initiated Sprint 4 by stabilizing evidence workflows, then implementing status lifecycle governance and approval foundations. I modeled transitions between PLANNED, EXECUTED, PENDING_VALIDATION, APPROVED, REJECTED, and ARCHIVED, and introduced StatusHistory tracking for auditability. I learned the value of explicit state machines in preventing inconsistent business states. I used domain modeling, transaction management, and test-driven validation of transition rules. A challenge was guaranteeing that invalid transitions could never silently pass through business logic. I solved this by centralizing transition checks in service methods and raising clear domain exceptions for invalid operations. My reflection is that governance logic is where business reliability is truly enforced.

\section*{Week 13}
This week I implemented the complete approval and rejection workflow for supervisors. I delivered endpoints such as /api/rapport-requests/{id}/approve, /api/rapport-requests/{id}/reject, pending queue retrieval, and status history consultation. On the Angular side, I built approval screens with rejection forms and feedback loops to technicians. I learned that quality assurance features must combine technical correctness with communication clarity. I used Spring transactional services, @PreAuthorize role protection, Angular reactive forms, and notification integration. The challenge was ensuring that rejections were constructive and traceable instead of becoming simple status changes. I solved this by persisting explicit rejection reason/feedback data and linking it to notification flows. My reflection is that effective systems support decision-making, not just state updates.

\section*{Week 14}
I launched Sprint 5 by implementing Excel report generation and configurable export definitions. I built report services using Apache POI, data extraction based on filters, dynamic column mapping, and downloadable XLSX generation for coordinators. I learned how reporting quality depends on both data modeling and export configuration discipline. I used Java IO handling, workbook styling/population, and repository query specification skills. The challenge was transforming heterogeneous intervention data into consistent client-ready spreadsheets. I solved this by formalizing ConfigExport mappings and building field extraction logic that gracefully handles missing or optional values. This week confirmed that reporting is where technical implementation becomes directly visible to business stakeholders.

\section*{Week 15}
In the final week, I concentrated on optimization, testing, and deployment readiness. I improved query performance with eager fetch strategies and caching, set up archival logic for older interventions, reviewed monitoring metrics, and validated results through unit, integration, and performance tests. I learned how final project quality comes from hardening and verification, not only feature completion. I used JPA optimization, cache configuration, scheduled task design, CI-oriented testing, and end-to-end validation skills. The main challenge was balancing optimization efforts with functional stability at the end of the timeline. I solved this by prioritizing bottlenecks with measurable impact, then verifying no regression through targeted test runs. My final reflection is that this internship significantly strengthened my full-stack engineering maturity: I progressed from requirement analysis to secure implementation, from mobile resilience to reporting automation, and from feature delivery to production-level quality thinking.

\section*{Overall Reflection}
Throughout this internship, I developed a strong understanding of how to build an end-to-end professional system that connects business process governance, multi-platform development, and operational reliability. I improved my technical depth in Spring Boot, Angular, Flutter, PostgreSQL, API design, notification workflows, and report automation, while also improving organization, prioritization, and communication in a sprint-driven process. The most important lesson I take from this experience is that engineering value is created when technical decisions remain continuously aligned with field reality and user needs.